下面是正确答案。

\section{第 1…6 章的小测验}

\begin{enumerate}
\item 2；顺便一提，对于带有基类的类，你可以使用继承构造函数（inheriting constructor），它使用基类名而非派生类名来声明构造函数。
\item 2
\item 2
\item 3
\item 2
\item 2、3
\item 1；\texttt{auto} 遵循模板类型推导（template type deduction）的规则，因此引用性和 const 限定不会被保留。
\item 1、2、3 —— 三个答案都正确
\item 3 —— 因为我们使用 \texttt{auto elem : vec}，\texttt{elem} 是向量中某个元素的副本，所以即便修改它，向量中的值也不会受影响。
\item 2、3 —— 右值引用（\texttt{42}）可以绑定到 \texttt{const} 引用或普通值。
\end{enumerate}

\section{最后的小测验}

\begin{enumerate}
\item 2
\item 2
\item 1
\item 1
\item 3
\item 2
\item 2
\item 1
\item 2；聚合体（aggregate）与类内成员初始化（in-class member initialization）自 C++14 起可用。
\item 2；（移动构造函数会阻止隐式拷贝构造函数的生成，参见代码 @C++ Insight\footnote{\url{https://cppinsights.io/s/9a1daa06}}）
\item 2；正确的元素类型是 \texttt{std::pair<const std::string, int>}，因此每次循环迭代都会产生一次拷贝；参见"推导"一章。
\item 3；这是拷贝列表初始化（copy list initialization）的特殊规则，它会生成 \texttt{initializer\_list}。
\item 1；\texttt{opt1} 需要一个临时对象。
\item 3
\item 1；\texttt{initializer\_list} 要求类型可拷贝；参见"非规则数据成员"一章中"某些不便——不可拷贝类型"小节。
\end{enumerate}

\section{第一个编程题的解答（NSDMI）}

第一个编程题的解答。在 Compiler Explorer 上运行

\begin{codebox}[cpp]
struct Point {
double x { 1.0 };
double y { 2.0 };
};
\end{codebox}

如你所见，该解使用 NSDMI 将 \texttt{x} 和 \texttt{y} 初始化为所需的值。

\section{第二个编程题的解答（NSDMI）}

第二个编程题的解答。在 Compiler Explorer 上运行

\begin{codebox}[cpp]
constexpr unsigned int DEFAULT_CATEGORY = 4;
constexpr unsigned int DEFAULT_FLAGS = 0x0a;
struct SalesRecord {
std::string name_ {"empty"};
double price_ { 1.0 };
unsigned int category_ : 4 { DEFAULT_CATEGORY };
unsigned int flags_ : 4 { DEFAULT_FLAGS};
};
\end{codebox}

该解将数据成员初始化为所需的值，包括自 C++20 起支持的位字段（bit field）。

\section{第三个编程题的解答（inline）}

计数类型（Counted Type）题目的解答。在 Wandbox 上运行

\begin{codebox}[cpp]
struct CountedType {
static inline int instanceCounter = 0;
static inline int maxInstanceCounter = 0;
// simple counting... only ctor and dtor implemented...
CountedType() { ++instanceCounter; ++maxInstanceCounter; }
~CountedType() { --instanceCounter; }
CountedType(const CountedType&) { ++instanceCounter; ++maxInstanceCounter; }
};
\end{codebox}

该解实现了默认构造函数、拷贝构造函数和析构函数。由于我们想要知道实例的最大数量，这个变量在析构函数中不会被递减。

\section{第四个编程题的解答（修正代码）}

第四题的解答。在 Compiler Explorer 上运行

\begin{codebox}[cpp]
struct SalesRec {
std::string name_;
double price_{}; // << make it 0 by default!
};
void addPromo(std::vector<SalesRec>& sales, double discount) {
for (auto& elem : sales)
// << reference
elem.price_ = (1.0-discount)*elem.price_;
}
double computeTotal(const std::vector<SalesRec>& sales) {
double sum{}; // << set it to 0 at start
for (const auto& elem : sales) // << don't copy elements
sum += elem.price_;
return sum;
}
\end{codebox}

该解在 4 处使用了正确的语法。它强制让 \texttt{price\_} 和 \texttt{sum} 在开始时被正确初始化为 \texttt{0.0}，并在循环迭代中使用了恰当的语义。
